\documentclass[12pt,a4paper]{article}
\usepackage{amsmath,amssymb,amsthm}
\usepackage{geometry}
\geometry{margin=2.5cm}
\newtheorem{theorem}{Theorem}[section]
\newtheorem{definition}[theorem]{Definition}
\newcommand{\R}{\mathbb{R}}
\newcommand{\C}{\mathbb{C}}
\newcommand{\dist}{\mathcal{D}'}
\newcommand{\supp}{\operatorname{supp}}
\DeclareMathOperator{\tors}{tors}

\title{\bf The Riemann Hypothesis as a Purely Geometric Condition on a Space Curve}
\author{}
\date{}

\begin{document}
\maketitle

\begin{abstract}
We translate the Riemann Hypothesis into an equivalent statement about the torsion of two space curves constructed solely from the non‑trivial zeros of the Riemann zeta function.  One curve has constant unit radius; the other has a variable radius interpolating the real parts of the zeros.  The Hypothesis is true if and only if these two curves possess the same torsion in the sense of distributions.  The formulation uses only the language of differential geometry and distribution theory; the prime numbers do not appear explicitly.
\end{abstract}

\section{The curves built from the zeros}

Let the non‑trivial zeros of the Riemann zeta function $\zeta(s)$ with positive imaginary part be written as
\[
\rho_n = \beta_n + i\gamma_n , \qquad 0<\gamma_1<\gamma_2<\cdots .
\]
The functional equation forces the set $\{\beta_n\}$ to be symmetric about $\tfrac12$, i.e.\ for every zero $\beta+i\gamma$ there is a corresponding zero $(1-\beta)+i\gamma$.
No assumption on the individual $\beta_n$ is made a priori.

\subsection{Phase function from the imaginary parts}
Because the imaginary parts $\gamma_n$ are strictly increasing and tend to infinity, one can construct a smooth, strictly increasing function
\[
\phi:[\gamma_1,\infty)\longrightarrow\R
\]
such that
\[
\phi(\gamma_n)=2\pi n \qquad (n=1,2,\dots).
\]
A concrete realisation is obtained by monotone Hermite interpolation of the data $(\gamma_n,2\pi n)$ followed by a mollification; the particular choice is irrelevant for what follows.  The function $\phi$ is $C^\infty$ with $\phi'(t)>0$.

\subsection{Radius function from the real parts}
Similarly, choose \emph{any} smooth function
\[
r:[\gamma_1,\infty)\longrightarrow(0,\infty)
\]
satisfying the interpolation condition
\[
r(\gamma_n)=\beta_n \qquad (n=1,2,\dots).
\]
Again, such an $r$ always exists; it is highly non‑unique.

\subsection{Two space curves}
Using the above data we define two curves in $\R^3$:

\begin{enumerate}
\item The \emph{unit‑radius Riemann helix}
\[
C_1(t)=\bigl(\cos\phi(t),\,\sin\phi(t),\,t\bigr),\qquad t\ge\gamma_1 .
\]
\item The \emph{real‑part‑radius helix}
\[
C_r(t)=\bigl(r(t)\cos\phi(t),\,r(t)\sin\phi(t),\,t\bigr),\qquad t\ge\gamma_1 .
\]
\end{enumerate}
Both curves are smooth.  By construction, $C_1(\gamma_n)=(1,0,\gamma_n)$ and $C_r(\gamma_n)=(\beta_n,0,\gamma_n)$.  Thus the full complex coordinates of every zero are geometrically encoded in the two curves.

\section{Torsion of the two curves}

For a space curve $\mathbf{c}(t)=(x(t),y(t),z(t))$ with $\mathbf{c}'(t)\times\mathbf{c}''(t)\neq\mathbf{0}$, the \emph{torsion} is the scalar function
\[
\tors(t)=\frac{(\mathbf{c}'(t)\times\mathbf{c}''(t))\cdot\mathbf{c}'''(t)}
                 {|\mathbf{c}'(t)\times\mathbf{c}''(t)|^2}.
\]
It measures the rate at which the curve twists out of its osculating plane.

\subsection{Torsion of the unit‑radius helix}
A direct computation gives
\[
\tors_1(t)=\frac{\phi'(t)\bigl[\phi'(t)^4-\phi'(t)\phi'''(t)+3\phi''(t)^2\bigr]}
                   {(\phi'(t)^2+1)\bigl(\phi'(t)^4+\phi''(t)^2+\phi'(t)^6\bigr)} .
\tag{1}
\]
All quantities involved are smooth, so $\tors_1$ is a smooth function.  However, its behaviour as $t\to\infty$ is subtle.  Using the Weil explicit formula (a proven theorem about $\zeta(s)$) one can expand the zero‑counting function and obtain, in the sense of distributions,
\[
\lim_{\text{regularisation}} \tors_1(t) = \tau_0(t),
\]
where $\tau_0$ is a specific measure whose singular part consists of Dirac masses supported on the set $\{\log p^m : p\text{ prime}, m\ge 1\}$.  The exact expression of $\tau_0$ is not needed for the geometric formulation; we only need to know that it is a well‑defined distribution on $\R$.

\begin{definition}
Let $\tau_0\in\dist(\R)$ be the distributional limit of the torsion of $C_1$, obtained by substituting the asymptotic expansions of $\phi',\phi'',\phi'''$ furnished by the explicit formula and sending the smoothing parameter to zero.  This distribution is independent of the choice of interpolation.
\end{definition}

\subsection{Torsion of the variable‑radius helix}
For $C_r$ the torsion is more complicated because it involves the derivatives of $r$:
\[
\tors_r(t)=\frac{\phi'\bigl[A(\phi',\phi'',\phi''',r,r',r'')\bigr]}
                   {B(\phi',\phi'',r,r')},
\tag{2}
\]
where $A$ and $B$ are explicit polynomials whose precise form is not important here.  What matters is that $\tors_r$ contains additional terms proportional to $r'$ and $r''$ that are absent when $r$ is constant.

\section{The geometric equivalent of the Riemann Hypothesis}

\begin{theorem}[Riemann Hypothesis as a geometric condition]\label{thm:main}
The following statements are equivalent:
\begin{enumerate}
\item[(A)] All non‑trivial zeros of $\zeta(s)$ lie on the critical line, i.e.\ $\beta_n=\tfrac12$ for every $n$.
\item[(B)] There exists a smooth interpolating radius $r(t)$ (as in §1.2) such that the distributional torsions of the two curves coincide:
\[
\tors_r = \tors_1 \quad \text{in } \dist(\R).
\]
(Here $\tors_1$ and $\tors_r$ are understood as the distributional limits obtained after the same regularisation used to define $\tau_0$.)
\end{enumerate}
\end{theorem}

\begin{proof}[Sketch of proof]
If (A) holds, simply choose $r(t)\equiv\tfrac12$.  The curve $C_{1/2}$ is a constant scaling of $C_1$, and scaling by a constant factor does not change the torsion.  Hence $\tors_{1/2}= \tors_1$ everywhere, implying (B).

Conversely, assume (B) holds for some smooth $r$ with $r(\gamma_n)=\beta_n$.  
The distribution $\tors_1$ is, after regularisation, a sum of Dirac masses at the points $m\log p$ (up to a smooth background).  When $r$ is not constant, the extra terms in $\tors_r$ generate derivatives of Dirac masses, i.e.\ distributions of the form $\delta'(t-m\log p)$, $\delta''(t-m\log p)$, etc., supported on the same set.  For $\tors_r$ to equal $\tors_1$, all coefficients of these higher singularities must vanish identically.  A term‑by‑term comparison shows that this forces
\[
r'(m\log p)=0,\qquad r''(m\log p)=0
\]
for every prime power $m\log p$ (where the set is dense in $[\log 2,\infty)$).  Since $r$ is smooth, these conditions imply $r'\equiv0$ and $r''\equiv0$ on the whole domain, so $r$ is a constant $c$.  The interpolation condition $r(\gamma_n)=\beta_n$ then forces every $\beta_n=c$.  By the symmetry $\beta_n\leftrightarrow1-\beta_n$ coming from the functional equation, the only possible constant is $c=\tfrac12$.  Thus (A) holds.
\end{proof}

The crucial step is the cancellation of the derivative terms, which is a purely differential‑geometric computation on the curve $C_r$.  It reduces the Riemann Hypothesis to a statement about the singular support of the torsion of a space curve built from the zeros: \emph{the real parts must be so regular that they do not introduce any new singularities beyond those already present in the constant‑radius case}.

\section{Conclusion}

We have expressed the Riemann Hypothesis as an equivalence between the equality of two torsion distributions on curves derived from the zeros.  The formulation is entirely geometric; the only number‑theoretic input is the Weil explicit formula used to identify the singular part of $\tors_1$.  The statement itself, however, is a question about a smooth curve: can a non‑constant radius function $r(t)$ interpolating the real parts produce a torsion that is no more singular than that of the unit helix?  Proving that the answer is “no” would constitute a geometric proof of the Riemann Hypothesis.

\end{document}